%% ccs-xml-test — CCS concepts pasted from the ACM CCS tool (twin of ccs-xml-test.typ).
%% LaTeX typesets the \ccsdesc lines (CCSXML is a comment environment); the Typst twin receives the equivalent <ccs2012> XML.
%% The set exercises grouping of an interleaved repeated area, the 500 (bold) / 300 (italic) / default-100 (roman) significance styles, and a trailing area-only repeat (\ccsdesc[100]{Networks}) whose only visible effect is acmart's @concepts counter quirk: the list ends "; " instead of ".".
\documentclass[acmsmall]{acmart}

\acmJournal{JACM}
\acmVolume{37}
\acmNumber{4}
\acmArticle{111}
\acmYear{2018}
\acmMonth{8}
\acmDOI{XXXXXXX.XXXXXXX}
\setcopyright{acmlicensed}
\copyrightyear{2018}

\begin{document}

\title{Classified Concepts}

\author{Grace Hopper}
\email{hopper@example.org}
\affiliation{%
  \institution{Compiler Institute}
  \city{Arlington}
  \country{USA}}

\begin{abstract}
  A short note whose only purpose is to typeset a CCS Concepts line parsed from the ACM CCS tool's output.
\end{abstract}

\begin{CCSXML}
<ccs2012>
 <concept>
  <concept_id>10010520.10010553.10010562</concept_id>
  <concept_desc>Computer systems organization~Embedded systems</concept_desc>
  <concept_significance>500</concept_significance>
 </concept>
 <concept>
  <concept_id>10003033.10003083.10003095</concept_id>
  <concept_desc>Networks~Network reliability</concept_desc>
  <concept_significance>300</concept_significance>
 </concept>
 <concept>
  <concept_id>10010520.10010553.10010554</concept_id>
  <concept_desc>Computer systems organization~Robotics</concept_desc>
  <concept_significance>100</concept_significance>
 </concept>
 <concept>
  <concept_id>10003033</concept_id>
  <concept_desc>Networks</concept_desc>
  <concept_significance>100</concept_significance>
 </concept>
</ccs2012>
\end{CCSXML}

\ccsdesc[500]{Computer systems organization~Embedded systems}
\ccsdesc[300]{Networks~Network reliability}
\ccsdesc{Computer systems organization~Robotics}
\ccsdesc[100]{Networks}

\keywords{classification, concepts}

\maketitle

\section{Body}
The CCS line above groups the interleaved repeated area, styles the specifics by significance, and ends without a period because of the area-only concept.

\end{document}
